\chapter{검증: 우연이 아님을 보이는 법}

``1.03\%p 더 번다''는 숫자 하나는 아무것도 증명하지 않는다. 시드가 좋았을 수도, 특정 해가 좋았을 수도, 수십 가지를 시도하다 하나가 걸렸을 수도 있다. 이 장은 그 가능성을 하나씩 지우는 도구들이다.

\section{사전 등록: 결과를 보기 전에 기준을 정한다}
11단계부터 모든 실험은 \textbf{판정 기준을 코드 맨 위 설명에 먼저 적고} 실행했다. 예: ``가장 단순한 방식보다 4개 연도 중 3개 이상에서 선별 가치 0.1\%p 또는 AUC 0.002 넘게 나으면 채택, 없으면 단순한 방식.''

\begin{keybox}[왜 중요한가]
결과를 보고 기준을 정하면, 어떤 결과든 ``성공''으로 만들 수 있다(garden of forking paths). 기준을 먼저 적어 두면 실패도 실패로 기록된다. 이 프로젝트의 실패 목록(연구 노트 2절)이 긴 것은 이 원칙의 결과다.
\Intuit{과녁을 먼저 그리고 쏴야 한다. 쏜 뒤에 화살 주위에 과녁을 그리면 백발백중이다.}
\end{keybox}

test 데이터에 대해서는 한 단계 더 나갔다. 평가 규칙, 모델, 가설, 지지 기준을 \ft{docs/test_protocol.md}에 적고 git 태그 \ft{test-protocol-v1}을 달았다. test를 본 뒤 바꿔야 하면 문서를 고치지 않고 ``이탈 기록''에 남긴다.

\section{위약 검정과 시드}
\begin{itemize}[leftmargin=*]
\item \textbf{라벨 섞기(placebo)}: 학습 라벨을 무작위로 섞어 모델을 만들면, 선별 가치는 0 근처여야 한다. 실제로 0 근처였다. 평가 코드 자체가 양수를 만들어 내는 버그가 없다는 확인이다.
\item \textbf{무작위 승인}: 같은 비율을 무작위로 승인한 포트폴리오를 기준선으로 뺀다(\S6.4).
\item \textbf{시드}: 모델 학습의 무작위성(행/열 표본추출, early stopping 분할)만 바꿔도 선별 가치가 최대 0.09\%p 움직인다. 효과가 이보다 작으면 의미가 없다(C16). 그래서 비교는 시드 3--10개 평균으로 한다.
\end{itemize}

\section{발행월 부트스트랩}
\textbf{질문.} 선별 가치 1.03\%p의 불확실성은 얼마인가?

같은 달 발행된 대출들은 같은 경기를 겪으므로 서로 독립이 아니다. 대출 단위로 다시 뽑으면 불확실성을 과소평가한다. 그래서 \textbf{발행월}을 단위로 다시 뽑는다(block bootstrap). 선별 가치는 월별 값 $f_m$의 평균이므로, 48개 월을 복원추출로 48번 뽑아 평균 내기를 1만 번 반복하고, 그 분포의 2.5\%, 97.5\% 분위를 95\% 구간으로 쓴다.

\begin{workout}[월별 차이 5개의 부트스트랩]
\Use{bootstrap(이 절)}
\Step{1} 예시: 두 모델의 월별 선별 가치 차이 $(1.2, 0.8, -0.2, 1.0, 0.6)$\%p. 평균 0.68\%p.
\Step{2} 5개 중 5개를 복원추출(예: 0.8, 0.8, 1.0, $-0.2$, 1.2)하고 평균을 낸다.
\Step{3} 1만 번 반복해 평균 1만 개를 얻는다. \Concept{bootstrap 분포}
\Step{4} 2.5\%와 97.5\% 분위: 0.24\%p와 1.04\%p.
\Ans{95\% 구간 (0.24, 1.04)\%p. 0을 포함하지 않으므로 ``차이가 있다''.}
\Verify{\ft{checks.py} 11번(시드 0).}
\Trap{두 모델을 따로 부트스트랩해 구간이 겹치는지 본다. 같은 달끼리 짝지은 \textbf{차이}를 부트스트랩해야 공통 변동이 지워져 훨씬 정밀하다(paired).}
\end{workout}

\section{여러 번 비교할 때: Holm 보정}
4개 모델을 기준과 비교하면 검정이 4번이다. 각각 5\% 수준이면 하나쯤은 우연히 ``유의''가 나올 수 있다. Holm 보정은 $p$값을 작은 순으로 정렬하고 $j$번째에 $(k-j+1)$을 곱한 뒤, 앞의 값보다 작아지지 않게 누적 최댓값을 취한다.

\begin{workout}[Holm 보정]
\Use{Holm(이 절)}
\Step{1} 예시: $k=3$, $p=(0.01, 0.02, 0.04)$.
\Step{2} $0.01\times3=0.03$, $0.02\times2=0.04$, $0.04\times1=0.04$.
\Step{3} 누적 최댓값: $(0.03, 0.04, 0.04)$.
\Ans{세 개 모두 0.05 미만이라 유의. Bonferroni(모두 \ensuremath{\times}3)였다면 0.04가 0.12가 되어 세 번째는 탈락.}
\Verify{\ft{checks.py} 12번.}
\end{workout}

\section{데이터 엿보기와 보류 구간}
2013--16년 평가 구간을 7단계부터 26단계까지 계속 다시 썼다. 규칙, 지표, 피처, 모델을 이 구간에서 고를수록 이 구간에 맞춰질 위험이 커진다. 지금까지 비교한 규칙, 지표, 모델 조합은 약 80개다.

\begin{itemize}[leftmargin=*]
\item 12단계 EDA는 2007--12년만 봤다.
\item 14단계에서 처음으로 2017년 발행분을 보류 구간으로 썼다(선별 가치 1.22\%p, 다만 종결 편향 있음). 이후 이 구간도 ``본 구간''이다.
\item 23--26단계 결과가 좋아도 test 프로토콜의 모델은 바꾸지 않았다.
\item 최종 판단은 한 번도 보지 않은 공식 test로만 한다.
\end{itemize}

\begin{warnbox}[test 해석 규칙(26단계에서 미리 정함)]
test에도 진행 중 대출이 없으므로 종결 편향이 생긴다. 발행연도별 만기 도달 비율부터 확인하고, 비율이 높은 연도만 기대수익 순위 vs 부도확률 순위 비교의 근거로 쓴다. 비율이 낮은 연도에서 부도확률 순위가 이기면 ``예상된 편향 방향''으로 보고하고, 그것만으로 C6를 기각하지 않는다.
\end{warnbox}

\section{주장 장부: 확신의 등급}
모든 주장은 \ft{docs/research_log.md}의 장부에 근거 파일, 상태(지지/잠정/기각), 남은 약점과 함께 적는다. 예:

\begin{table}[H]\centering\small
\begin{tabularx}{\linewidth}{@{}p{8mm}Lp{20mm}L@{}}\toprule
\# & 주장 & 상태 & 남은 약점\\\midrule
C6 & 기대수익 순위가 부도확률 순위보다 더 번다 & 지지(전체 표본) & 격차가 해마다 줄어듦\\
C6b & 부도확률 순위는 무작위와 비슷하다 & \textbf{기각} & 시드 5개로 보면 유의하게 양수\\
C7 & 가치는 등급 안 선별에서 나온다 & 지지 & 2013년은 배분이 주효과\\
C19 & AUC 순위는 돈 순위와 다르다 & 지지(탐색) & 수준은 극단적일 때만 문제\\
C23 & 종결 편향은 부도확률 순위를 유리하게 만든다 & 지지 & 2013--15년 36개월만\\\bottomrule
\end{tabularx}
\caption{주장 장부 발췌.}
\end{table}

\begin{career}
``결과가 우연이 아니라는 걸 어떻게 아나요?''에 대한 답이 이 장 전체다. 사전 등록, 위약 검정, 다중 시드, 군집 부트스트랩, 다중비교 보정, 보류 구간, 그리고 ``틀린 주장을 기각한 기록''. 특히 마지막이 신뢰를 만든다.
\end{career}
